\begin{tikzpicture}[>=stealth,font=\small]
\draw[gray!40,rounded corners] (0,0) rectangle (14,6.8);
\node at (7,6.3) {关于 $x$ 轴对称：非零输出会冲突，零输出反射后仍是自己};
\begin{scope}[shift={(3.1,3.2)},x=1.3cm,y=1cm]
\draw[->,gray] (-1.5,0)--(2.2,0) node[right]{$x$};
\draw[->,gray] (0,-1.6)--(0,1.8) node[above]{$y$};
\draw[orange!80!black,dashed] (1,-1.5)--(1,1.5);
\fill[blue] (1,1) circle(2pt) node[above left]{$(1,1)$};
\fill[red!70!black] (1,-1) circle(2pt) node[below left]{$(1,-1)$};
\node[align=center] at (.3,-2.3) {若同时保留这两个点\\同一输入 $1$ 就有两个输出};
\end{scope}
\begin{scope}[shift={(10.3,3.2)},x=1.3cm,y=1cm]
\draw[->,gray] (-1.7,0)--(2,0) node[right]{$x$};
\draw[->,gray] (0,-1.6)--(0,1.8) node[above]{$y$};
\draw[blue,very thick] (-1.6,0)--(1.7,0);
\draw[orange!80!black,dashed] (1,-1.4)--(1,1.5);
\fill[orange!80!black] (1,0) circle(2pt) node[above left]{$(1,0)$};
\node at (-.95,1.3){$f(x)\equiv0$};
\node[align=center] at (0,-2.3) {镜像点与原来的点重合\\同一输入仍只有输出 $0$};
\end{scope}
\end{tikzpicture}
